\section*{High-Throughput Textile Decontamination at the Center for High Performance Cleansing (CHPC)}

The University of Utah takes pride in innovation, efficiency, and supporting
students in every aspect of campus life.

To that end, the university has established:
\begin{quote}
\textbf{The Center for High Performance Cleansing (CHPC)}
\end{quote}

Officially, CHPC is ``a multidisciplinary facility for advanced hydrothermal
textile decontamination and surfactant dynamics research.''

To maximize both sustainability and student success, CHPC runs a single,
high-efficiency Decontamination Chamber (washer-dryer unit), and the university
would like to use it as effectively as possible.

Throughout the day, students, TAs, and faculty submit laundry ``experiments''
(loads). Each load has:
\begin{itemize}
    \item \(\texttt{s}\): the earliest time the load can start.\\
    This is when its owner and their laundry arrive at CHPC.
    \item \(\texttt{f}\): the time the load would finish if started at
    \(\texttt{s}\) and runs without interruption.
\end{itemize}

Once a load starts, it runs continuously until it finishes. No mid-cycle
interruptions, no pausing ``just for one more sock.''

Two loads conflict if their time intervals overlap in any way. Since CHPC is
committed to safe and reliable operation, the Decontamination Chamber can only
handle one load at a time.

The university wants to show that CHPC significantly increases ``experimental
throughput'' (and also clean clothes) by running as many loads as possible in a
day.

\medskip

\noindent
\textbf{Task.} What is the maximum number of non-conflicting loads that can be
scheduled in the Decontamination Chamber?
